\documentclass[letter,12pt]{article}
\usepackage[paperheight=27.94cm,paperwidth=21.59cm,bindingoffset=0in,
            left=3cm,right=2.0cm, top=3.5cm,bottom=2.5cm, 
            headheight=200pt, headsep=1.0\baselineskip]{geometry}
\usepackage{graphicx,lastpage}
\usepackage{upgreek}
\usepackage{censor}
\usepackage[spanish,es-tabla]{babel}
\usepackage{graphicx}
\usepackage{pdfpages}
\usepackage{tabularx}
\usepackage{adjustbox}
\usepackage{xcolor}
\usepackage{colortbl}
\usepackage{rotating}
\usepackage{multirow}
\usepackage{float}
\usepackage{enumitem}
\usepackage[utf8]{inputenc}
\renewcommand{\tablename}{Tabla}
\usepackage{fancyhdr}
\pagestyle{fancy}

\setlength{\emergencystretch}{3em}

\title{\Huge{Informe N° 3.}}
\author{Alex Nuñez  \& Camilo Pinto \& Lucas Gonzalez}
\date{Octubre de 2025}

\begin{document}

\maketitle
\tableofcontents
\newpage

\section{Descripción del sistema}
Se desarrollará un sistema que centraliza la gestión de solicitudes internas mediante tickets, permitiendo registrar incidentes/requerimientos con título, descripción, categoría, prioridad y adjuntos, para darles seguimiento trazable hasta su cierre.

El sistema incluirá autenticación y roles, creación y asignación de tickets, ciclo de vida (Nuevo → En progreso → Resuelto → Cerrado), comentarios e historial de cambios; además, listados con filtros (estado, prioridad, categoría, técnico, fechas) y un tablero para visualizar la carga. Generará un reporte resumido (resueltos/cerrados por categoría), enviará notificaciones en eventos, permitirá la reapertura si el problema persiste, además de la eliminación/anulación por el solicitante cuando corresponda.

El cliente principal estará basado en tecnologías web y no requerirán instalación de software adicional, será utilizado por todos los usuarios. Mientras que los servicios principales serán: Tickets, Auth y Notificaciones.

\section{Objetivos esperados}
\begin{enumerate}
    \item Gestionar tickets de ayuda (help desk): Permitir la creación, seguimiento, asignación, y cierre de tickets.
    \item Mejorar la comunicación y resolución de problemas: Facilitar la interacción entre usuarios y administradores, y la posible reapertura de tickets.
    \item Proporcionar un control y visibilidad del estado de los problemas: Ofrecer funcionalidades de listado, filtros, y reportes resumidos para el seguimiento de los tickets.
    \item Optimizar la atención al usuario: A través de notificaciones y la capacidad de eliminación de tickets por parte del solicitante.
\end{enumerate}
\section{Perfiles de usuario}
El sistema está dirigido a los colaboradores internos de la organización y define tres roles de usuario principales, cada uno con funcionalidades específicas:

\begin{enumerate}
    \item Usuario: Es el rol estándar para todos los colaboradores. Pueden crear nuevos tickets para reportar incidentes o hacer requerimientos, añadir comentarios, adjuntar archivos, dar seguimiento al estado de sus solicitudes, y reabrir o anular sus propios tickets.

    \item Coordinador/Técnico: Es el usuario encargado de resolver los tickets. Tiene la capacidad de visualizar el tablero con la carga de trabajo, tomar la asignación de tickets, cambiar su estado (e.g., a "En Progreso" o "Resuelto"), y comunicarse con los solicitantes a través de comentarios.

    \item Administrador: Posee el nivel de acceso más alto. Además de las funcionalidades de los otros roles, puede gestionar usuarios, categorías, y tiene acceso a todos los tickets y reportes para obtener una visión global del sistema, lo que le permite tomar decisiones estratégicas.
\end{enumerate}
\section{Requerimientos}
\begin{enumerate}[leftmargin=*]
    \item Autenticación y roles: login/logout, recuperación de contraseña, control de acceso por rol (Administrador/Coordinador/Usuario), bloqueo por intentos y expiración de sesión.
    \item Gestión de tickets (CRUD + ciclo de vida): crear/ver/actualizar/eliminación lógica con auditoría, adjuntos y comentarios; asignación a coordinadores; estados Nuevo $\rightarrow$ En progreso $\rightarrow$ Resuelto $\rightarrow$ Cerrado; reapertura desde Resuelto/Cerrado; registro de autor, fecha y cambios.
    \item Listado y filtros: listar y filtrar por estado, prioridad, categoría, asignado y fechas; paginación, ordenamiento y exportación CSV.
    \item Reporte resumen: gráfico de tickets Resueltos/Cerrados por categoría y período; indicadores de tiempo de resolución y tasa de reapertura; acceso según rol.
    \item Notificaciones: in-app y correo en creación, asignación, cambio de estado, comentarios y cierre; cola con reintentos y envío diferido si no hay conexión.
    \item Acceso y comunicación: cliente por línea de comandos (CLI); autenticación en terminal; interacción entre servicios vía bus externo.
    \item Arquitectura y protocolos: SOA con bus externo (\texttt{jrgiadach/soabus}) por TCP 5000; servicios desacoplados que publican/consumen mensajes; cliente CLI;\\
    \hspace*{1.5em}variables \texttt{BUS\_HOST=soabus} y \texttt{BUS\_PORT=5000}; contenedores Docker.
    \item Geolocalización y riesgos: geolocalización opcional con consentimiento (uso de permisos de navegador o ingreso manual con pin en mapa).
    \item Asignación automática de tickets: búsqueda de asignaciones de los operarios en base de datos, prioridad de asignación según actividad de operario y parámetros binarios; se implementará mediante el uso de FIFO y equilibradores de carga (borrador: ``asigna mediante reglas de negocio'' — categoría, prioridad, ubicación — y notifica al administrador y al usuario de asignación).
\end{enumerate}
\subsection{Requerimientos no funcionales}
\begin{enumerate}[leftmargin=*]
    \item Disponibilidad: servicio operativo con mantenciones planificadas.
    \item Seguridad: todo por conexiones seguras con cifrado vigente; hash seguro; autorización por rol; protección CSRF; límite de consultas; auditoría.
    \item Observabilidad: logs estructurados, métricas y trazas distribuidas.
    \item Cumplimiento y privacidad: ley chilena de datos personales; consentimiento para ubicación.
    \item Límites y restricciones: usuarios, tickets, tamaño de adjuntos y almacenamiento documentados.
    \item Riesgos y mitigación: copias de seguridad, reintentos de conexión.
\end{enumerate}
\subsection{Trazabilidad u objetivos}
\begin{enumerate}[leftmargin=*]
    \item RF Autenticación y roles $\rightarrow$ Objetivo de control de acceso.
    \item RF Gestión de tickets $\rightarrow$ Objetivo de gestión y trazabilidad.
    \item RF Asignación $\rightarrow$ Objetivo de resolución eficiente.
    \item RF Comentarios y adjuntos $\rightarrow$ Objetivo de colaboración.
    \item RF Listado y filtros $\rightarrow$ Objetivo de visibilidad del estado.
    \item RF Reporte resumen $\rightarrow$ Objetivo de control y decisiones.
    \item RF Notificaciones $\rightarrow$ Objetivo de atención oportuna.
    \item RF Acceso y comunicación $\rightarrow$ Objetivo de disponibilidad a usuarios.
    \item RF Geolocalización $\rightarrow$ Objetivo de análisis y despacho.
\end{enumerate}
\section{Modelo de datos y Persistencia}
Mecanismo de Persistencia: Para el almacenamiento de los datos se utilizará un sistema de gestión de bases de datos relacional (RDBMS). Esta elección se debe a la naturaleza estructurada de los datos, la necesidad de mantener la integridad referencial entre entidades como tickets, usuarios y comentarios, y las capacidades transaccionales que ofrece.

\begin{enumerate}[leftmargin=*]
    \item Tickets:
    \begin{itemize}[leftmargin=1.5em,noitemsep,topsep=0pt]
        \item ID (int pk, autoincrement).
        \item Título (varchar(120) not null).
        \item Descripción (text not null).
        \item categoria\_id (int fk $\rightarrow$ categorias.id, ON DELETE RESTRICT).
        \item Prioridad (enum: BAJA, MEDIA, ALTA, CRITICA).
        \item Estado (enum: NUEVO, EN\_PROGRESO, RESUELTO, CERRADO).
        \item Fecha creación (datetime default CURRENT\_TIMESTAMP).
        \item Fecha cierre (datetime null).
        \item solicitante\_id (int fk $\rightarrow$ usuarios.id, ON DELETE RESTRICT).
        \item ejecutivo\_id (int fk $\rightarrow$ usuarios.id, ON DELETE SET NULL).
        \item Índices: (estado, prioridad), categoria\_id, solicitante\_id, ejecutivo\_id.
    \end{itemize}

    \item Usuarios:
    \begin{itemize}[leftmargin=1.5em,noitemsep,topsep=0pt]
        \item ID (int pk, autoincrement).
        \item Nombre (varchar(120) not null).
        \item Email (varchar(120) unique not null).
        \item password\_hash (char(60) not null).
        \item Rol (enum: ADMIN, COORDINADOR, USUARIO).
        \item Fecha creación (datetime default CURRENT\_TIMESTAMP).
        \item Activo (tinyint(1) default 1).
    \end{itemize}

    \item Categorías:
    \begin{itemize}[leftmargin=1.5em,noitemsep,topsep=0pt]
        \item ID (int pk, autoincrement).
        \item Nombre (varchar(80) unique not null).
        \item Descripción (text null).
    \end{itemize}

    \item Comentarios:
    \begin{itemize}[leftmargin=1.5em,noitemsep,topsep=0pt]
        \item ID (int pk, autoincrement).
        \item Contenido (text not null).
        \item Fecha creación (datetime default CURRENT\_TIMESTAMP).
        \item ticket\_id (int fk $\rightarrow$ tickets.id, ON DELETE CASCADE).
        \item usuario\_id (int fk $\rightarrow$ usuarios.id, ON DELETE SET NULL).
        \item Índice: ticket\_id.
    \end{itemize}

    \item Adjuntos:
    \begin{itemize}[leftmargin=1.5em,noitemsep,topsep=0pt]
        \item ID (int pk, autoincrement).
        \item Nombre (varchar(120) not null).
        \item Ruta (varchar(255) not null).
        \item Tipo (varchar(100) not null).
        \item Tamaño (int unsigned null).
        \item Checksum (varchar(64) null).
        \item Fecha subida (datetime default CURRENT\_TIMESTAMP).
        \item ticket\_id (int fk $\rightarrow$ tickets.id, ON DELETE CASCADE).
        \item usuario\_id (int fk $\rightarrow$ usuarios.id, ON DELETE SET NULL).
        \item Índice: ticket\_id.
    \end{itemize}
\end{enumerate}

Conforme a los requisitos del proyecto, el sistema se diseñará bajo una Arquitectura Orientada a Servicios (SOA). Esta arquitectura desacopla las funcionalidades en servicios independientes y reutilizables, que se comunican a través de una red.  
Los componentes principales de la arquitectura son:
\begin{enumerate}[leftmargin=*]
    \item Componente Cliente: Aplicación por línea de comandos que invoca casos de uso 
y publica/consume mensajes vía bus. Sin UI web en esta entrega.

    \item Componentes de Servicio: Son los encargados de implementar la lógica de negocio. Se han definido tres servicios principales:  
\begin{itemize}
    \item Servicio de Autenticación (Auth): Responsable de gestionar la identidad de los usuarios: login, logout, gestión de sesiones, recuperación de contraseñas y validación de roles y permisos.
    \item Servicio de Tickets: Es el servicio principal que encapsula toda la lógica relacionada con la gestión de tickets: creación, actualización, asignación, comentarios, adjuntos, listados y reportes.
    \item Servicio de Notificaciones: Opera de forma asíncrona. Su función es gestionar el envío de correos electrónicos y notificaciones in-app, utilizando el bus externo para asegurar la entrega con reintentos incluso si los sistemas externos no están disponibles temporalmente.
\end{itemize}
\end{enumerate}
Se describen los contratos de mensajería por bus para cada servicio, indicando el requerimiento funcional (RF) que satisfacen.  

 \subsection*{Contratos por bus de servicios}
En esta entrega los servicios no exponen HTTP. Se comunican por un \textbf{bus TCP} externo
(\texttt{jrgiadach/soabus}, puerto 5000). El protocolo es \emph{JSON por línea} (JSONL): cada
mensaje es un JSON terminado en salto de línea. Hay dos patrones:
\textbf{request/response} (sincronía punto a punto) y \textbf{eventos} (publicación asíncrona).

\noindent\textbf{Convenciones}
\begin{itemize}
  \item Variables: \texttt{BUS\_HOST=soabus}, \texttt{BUS\_PORT=5000}.
  \item Campos: \texttt{type $\{$request|event$\}$}, \texttt{service}, \texttt{action} o \texttt{topic}, \texttt{payload}.
  \item Fiabilidad: el emisor reintenta ante error y el consumidor idempotiza por \texttt{id} de mensaje.
\end{itemize}

\paragraph{\texttt{Auth.authenticate (request)}} Autentica credenciales y retorna perfil.
\begin{verbatim}
{"type":"request","service":"auth","action":"authenticate",
 "payload":{"email":"usuario@acme.com","password":"****"}}
\end{verbatim}

\paragraph{\texttt{Tickets.get\_by\_user (request)}} Lista tickets del usuario.
\begin{verbatim}
{"type":"request","service":"tickets","action":"get_by_user",
 "payload":{"user_id":123}}
\end{verbatim}

\paragraph{\texttt{Ticket.created (event)}} Se emite al crear un ticket; lo consume Notificaciones.
\begin{verbatim}
{"type":"event","topic":"ticket.created",
 "payload":{"id":101,"categoria":"Soporte","usuario":123}}
\end{verbatim}

\subsection{Interfaz del Servicio de Notificaciones}
    
    Este servicio no expone una API pública para el cliente web, sino que se comunica a través del bus externo (jrgiadach/soabus) y aplica reintentos ante fallas externas. Otros servicios (como el de Tickets) publican eventos en el bus externo (e.g., {ticket.created}, {ticket.comment.added}) y el servicio de notificaciones los consume para realizar los envíos.
   \begin{enumerate}
    \item \textbf{Evento: ticket.created}
          \newline Descripción: Se activa al crear un ticket. El servicio envía una notificación por correo al solicitante.
          \newline RF Satisfecho: Notificaciones.
          \newline RNF Considerado: Disponibilidad (la cola con reintentos mitiga caídas del servidor de correo).
          
    \item \textbf{Evento: ticket.assigned}
          \newline Descripción: Se activa al asignar un ticket a un coordinador. El servicio notifica al usuario asignado.
          \newline RF Satisfecho: Notificaciones.
          
\end{enumerate}

\section{Procesos cliente}
Flujos de uso del sistema con evidencia.
\section{Procesos cliente (CLI)}
\subsection*{Crear ticket}
\begin{verbatim}
docker compose up -d soabus
docker compose run --rm -it app python -m Sistema_Ticket_Soa.main
# Login -> Crear ticket -> completar -> confirmar
\end{verbatim}

\subsection*{Asignar y comentar}
\begin{verbatim}
# Menú -> Asignar ticket -> seleccionar técnico
# Menú -> Agregar comentario -> escribir texto
\end{verbatim}

\subsection*{Cerrar y reabrir}
\begin{verbatim}
# Menú -> Cambiar estado a Resuelto/Cerrado
# Menú -> Reabrir si corresponde
\end{verbatim}

\subsection*{Crear ticket}
\begin{enumerate}
  \item Usuario inicia sesión.
  \item Abre ``Nuevo ticket'', completa y envía.
  \item Sistema confirma con ID y estado \textit{Abierto}.
\end{enumerate}
\begin{figure}[h]
\centering
\includegraphics[width=0.85\linewidth]{evidencias/proceso_crear_ticket.png}
\caption{Flujo crear ticket}
\end{figure}

\subsection*{Asignar y comentar}
Pasos para asignar responsable y agregar comentarios con captura de pantalla.
\begin{figure}[h]
\centering
\includegraphics[width=0.85\linewidth]{evidencias/proceso_asignar_comentar.png}
\caption{Flujo asignar y comentar}
\end{figure}

\subsection*{Cerrar y reabrir}
Criterios de cierre, permisos para reabrir y evidencia.
\begin{figure}[h]
\centering
\includegraphics[width=0.85\linewidth]{evidencias/proceso_cerrar_reabrir.png}
\caption{Flujo cerrar y reabrir}
\end{figure}
% \begin{figure}[h]
% \centering
% \includegraphics[width=0.85\linewidth]{evidencias/proceso_crear_ticket.png}
% \caption{Flujo crear ticket}
% \end{figure}

\subsection*{Asignar y comentar}
Pasos para asignar responsable y agregar comentarios con captura de pantalla.
% Temporarily commented out missing image
% \begin{figure}[h]
% \centering
% \includegraphics[width=0.85\linewidth]{evidencias/proceso_asignar_comentar.png}
% \caption{Flujo asignar y comentar}
% \end{figure}

\subsection*{Cerrar y reabrir}
Criterios de cierre, permisos para reabrir y evidencia.
% \begin{figure}[h]
% \centering
% \includegraphics[width=0.85\linewidth]{evidencias/proceso_Cambiar_contrasena.png}
% \caption{Flujo proceso Cambiar contrasena}
% \end{figure}


\section{Servicios implementados}
Estado por módulo, contratos y decisiones.
\subsection*{Resumen}
\begin{tabular}{|l|l|l|}
\hline
\textbf{Módulo} & \textbf{Descripción} & \textbf{Estado} \\
\hline
Autenticación & Login JWT, roles & Implementado/Parcial \\
Tickets & CRUD, asignación, comentarios, estados & Implementado/Parcial \\
Notificaciones & Email/In-app & Implementado/Parcial \\
\hline
\end{tabular}

\subsection*{Contratos de API (clave)}
\begin{tabular}{|l|l|l|}
\hline
\textbf{Método} & \textbf{Ruta} & \textbf{Descripción} \\
\hline
POST & /api/v1/auth/login & Devuelve token y perfil \\
GET  & /api/v1/tickets & Lista con filtros \\
POST & /api/v1/tickets & Crea ticket \\
GET  & /api/v1/tickets/\{id\} & Obtiene ticket \\
PUT  & /api/v1/tickets/\{id\} & Edita ticket \\
\hline
\end{tabular}



\section{Anexo}
\subsection*{Contratos por bus de servicios}
Los servicios intercambian mensajes JSON por línea (JSONL) en el bus TCP (\texttt{jrgiadach/soabus}).

\paragraph{\texttt{Auth.authenticate (request)}}
\begin{verbatim}
{"type":"request","service":"auth","action":"authenticate",
 "payload":{"email":"usuario@acme.com","password":"****"}}
\end{verbatim}

\paragraph{\texttt{Tickets.get\_by\_user (request)}}
\begin{verbatim}
{"type":"request","service":"tickets","action":"get_by_user",
 "payload":{"user_id":123}}
\end{verbatim}

\paragraph{\texttt{Ticket.created (event)}}
\begin{verbatim}
{"type":"event","topic":"ticket.created",
 "payload":{"id":101,"categoria":"Soporte","usuario":123}}
\end{verbatim}



\section{Comentarios}

Se decidió implementar para esta entrega un cliente por línea de comandos en vez de una interfaz web. La razón: bajar complejidad del frontend, menor tiempo requerido para su implementación y validar primero los servicios núcleo (Auth, Tickets, Notificaciones) y el flujo de negocio. El CLI permite crear, listar, asignar, comentar, cerrar y reabrir tickets usando el bus externo por TCP 5000 (jrgiadach/soabus), manteniendo trazabilidad y roles. La UI web queda como trabajo posterior, reutilizando los mismos contratos y agregando pruebas de usabilidad. Este recorte no afecta los RF clave ni los RNF de seguridad, observabilidad y disponibilidad planificados.

\end{document}